\section{Method}
\label{sec:method}

\subsection{Background: SVGD as a particle flow}

Let $p(x)$ denote a target density on $\mathbb{R}^d$ and let
$\{x_i^t\}_{i=1}^N$ be a set of particles at iteration $t$.  SVGD evolves particles by following a deterministic
interacting-particle flow whose velocity field is constrained to RKHS \cite{liu2016svgd}. A first-order discretization of
this flow is the update
\begin{equation}
  x_i^{t+1} = x_i^t + h\,\phi(x_i^t),
  \label{eq:svgd-update}
\end{equation}
where $h>0$ is a step size and the SVGD velocity field is
\begin{equation}
  \phi(x) = \frac{1}{N}\sum_{j=1}^N
  \Big[ k(x_j, x)\,\nabla_{x_j}\log p(x_j) + \nabla_{x_j} k(x_j, x) \Big].
  \label{eq:svgd-velocity}
\end{equation}
The first term in \cref{eq:svgd-velocity} is an \emph{attractive} drift that
pushes particles toward high-density regions via the score $\nabla \log p$. The second term is a \emph{repulsive} interaction that discourages particle collapse.
The multirate constructions below are kernel-agnostic: they apply once a
kernel-induced attractive field and repulsive field have been specified.

\subsection{Drift--repulsion splitting and multirate integration}

The central numerical observation motivating this work is that the drift and
repulsion components of \cref{eq:svgd-velocity} can impose different stability
constraints and can naturally evolve on different time scales. We therefore
write the SVGD flow as a split system
\begin{equation}
  \dot{x} = f_{\text{rep}}(x) + f_{\text{drift}}(x),
  \label{eq:svgd-splitting}
\end{equation}
where, for particle $i$,
\begin{align}
  f_{\text{drift}}(x_i)
  &:= \frac{1}{N}\sum_{j=1}^N k(x_j, x_i)\,\nabla_{x_j}\log p(x_j),
  \label{eq:svgd-drift-term}\\
  f_{\text{rep}}(x_i)
  &:= \frac{1}{N}\sum_{j=1}^N \nabla_{x_j} k(x_j, x_i).
  \label{eq:svgd-rep-term}
\end{align}
We next define multirate discretizations that apply different substepping
strategies to $f_{\text{drift}}$ and $f_{\text{rep}}$ while keeping a common
macro-step size $h$.

\paragraph{Strang-split SVGD}

Define the forward Euler update that advance the split system by a step size
$\tau$:
\begin{equation*}
  \Psi_{\text{rep}}^{\tau}(x) = x + \tau f_{\text{rep}}(x),
  \qquad
  \Psi_{\text{drift}}^{\tau}(x) = x + \tau f_{\text{drift}}(x).
\end{equation*}
\emph{Strang
splitting} \cite{strang1968difference} is a classical second-order method that
constructs a macro-step of size $h$ by composing such
substeps. In our setting, one macro step takes the form
\begin{equation}
  x^{t+1}
  = \Psi_{\text{rep}}^{h/2}\circ\Psi_{\text{drift}}^{h}\circ\Psi_{\text{rep}}^{h/2}(x^{t}),
  \label{eq:strang-split-svgd}
\end{equation}
which alternates half a repulsion step, a full drift step, and another half
repulsion step. In practice, the half-steps can be implemented as
$\frac{h}{2m}$ smaller substeps for stability.

\paragraph{Fixed multirate SVGD (MR-SVGD)}
Fixed multirate integrators provide a principled way to allocate different step
sizes to fast and slow components of a split system
\cite{Sandu_2016_MR-GARK,sarshar2019mrgark,sarshar2021implicitgark}. In our SVGD
setting, however, higher-order multirate schemes are often not cost-effective:
their additional stages translate into multiple evaluations of the drift and
repulsion fields per macro step, and each such evaluation is expensive due to
gradient computations and, in particular, $O(N^2)$ kernel interactions.
Moreover, practical accuracy is also limited by finite-particle effects and by
choices such as bandwidth selection and early stopping, which can mute the
benefits of increasing time-integration order at a fixed computational budget.
We therefore utilize a first-order multirate construction using Euler updates that preserves a macro step of size $h$ for the drift while resolving the
repulsion with $m$ substeps:
\begin{equation}
  x^{t+1}
  = \Psi_{\text{drift}}^{h}\circ\Big(\Psi_{\text{rep}}^{h/m}\Big)^{m}(x^{t}).
  \label{eq:mr-svgd}
\end{equation}
This retains a coarse drift step while stabilizing the repulsive interaction
when particles crowd, at the cost of additional kernel evaluations. The
corresponding implementation-level procedure is summarized in
\cref{alg:mr_svgd}.

\begin{algorithm}[t]
\caption{Fixed multirate SVGD (MR-SVGD)}
\label{alg:mr_svgd}
\begin{algorithmic}[1]
\Require Particles $\{x_i\}_{i=1}^N$, step size $h$, repulsion substeps $m$, kernel $k$
\State $x \leftarrow \{x_i\}$
\For{$s=1$ to $m$}
  \State compute repulsion $f_{\text{rep}}(x)$ using \cref{eq:svgd-rep-term}
  \State $x \leftarrow x + (h/m)\,f_{\text{rep}}(x)$
\EndFor
\State compute drift $f_{\text{drift}}(x)$ using \cref{eq:svgd-drift-term}
\State $x \leftarrow x + h\,f_{\text{drift}}(x)$
\State \Return $x$
\end{algorithmic}
\end{algorithm}

\paragraph{Adaptive multirate SVGD (Adapt-MR-SVGD)}
Although fixed-$m$ schedules are simple and robust, they impose a uniform drift
resolution along the entire trajectory. As a result, they may spend unnecessary
computational effort in benign regions while providing insufficient resolution
in locally stiff regimes. To address this imbalance, we choose the number of
drift microsteps $m$ adaptively at each macro step, while retaining one
repulsion update per macro step. This follows the standard local-error-control
perspective used in ODE
integrators, where local error estimates are used to adjust resolution
\cite[Ch. II]{Hairer_book_I}. After the repulsion update, denote the state by
$x=\Psi_{\text{rep}}^{h}(x^{t})$. We estimate the local
drift-discretization error by comparing one drift step of size
$h$ with two consecutive drift half-steps of size $h/2$,
\begin{align}
  x_{\text{full}} &= x + h\,f_{\text{drift}}(x), \\
  x_{\text{half}} &= x + \tfrac{h}{2}\,f_{\text{drift}}(x), \\
  \tilde{x}_{\text{full}}  &= x_{\text{half}} + \tfrac{h}{2}\,f_{\text{drift}}(x_{\text{half}}).
\end{align}
We then form the relative discrepancy
\begin{equation}
  \rho
  = \frac{\|\tilde{x}_{\text{full}} - x_{\text{full}}\|_2}{\|x\|_2 + \epsilon},
  \label{eq:adapt-mr-error}
\end{equation}
which serves as a dimensionless local error indicator for the drift update. $\epsilon>0$ is a small constant to avoid division by zero.
For a first-order drift integrator, this local error behaves
as \cite[Ch. II.4]{Hairer_book_I}:
\[
\rho(h) = C h^2 + \mathcal{O}(h^3),
\]
where $C$ is the principal error constant dependent on $f$ and its derivatives. If the drift step is replaced by $m$ microsteps of size $h/m$, then
\[
\rho_m \approx C(h/m)^2 \approx \rho/m^2.
\]
Imposing the target local error tolerance $\rho_m \approx \text{Tol}$ gives
$m \approx \sqrt{\rho/\text{Tol}}$. We therefore use the clipped integer
controller
\begin{equation}
  m \leftarrow \mathrm{clip}\Big(\Big\lceil \sqrt{\rho/\text{Tol}}\Big\rceil,
  m_{\min}, m_{\max}\Big),
  \label{eq:adapt-mr-m}
\end{equation}
which increases $m$ when $\rho>\text{Tol}$ and decreases it when
$\rho\le\text{Tol}$, while enforcing the robustness bounds
$m_{\min}\le m\le m_{\max}$.
The macro-step update is then
\begin{equation}
  x^{t+1}
  = \Big(\Psi_{\text{drift}}^{h/m}\Big)^m\circ\Psi_{\text{rep}}^{h}(x^t),
  \qquad m = m(x^t,h).
  \label{eq:adapt-mr-step}
\end{equation}
This
strategy increases drift resolution only when the local error indicator suggests
stiffness or rapid variation in the drift component.  \Cref{alg:adapt_mr_svgd} provides the complete algorithm for this method.


\begin{remark}
In practice, if $m\le 2$ we can reuse the already computed, locally-higher-order solution $\tilde{x}_{\text{full}} $ to avoid extra drift evaluations.
\end{remark}

\begin{algorithm}[t]
\caption{Adaptive multirate SVGD (Adapt-MR-SVGD)}
\label{alg:adapt_mr_svgd}
\begin{algorithmic}[1]
\Require Particles $\{x_i\}_{i=1}^N$, step size $h$, bounds $m_{\min}, m_{\max}$, tolerance $\text{Tol}$
\State $x \leftarrow \{x_i\}$
\State compute repulsion $f_{\text{rep}}(x)$ using \cref{eq:svgd-rep-term}
\State $x \leftarrow x + h\,f_{\text{rep}}(x)$
\State compute drift $f_{\text{drift}}(x)$ using \cref{eq:svgd-drift-term}
\State $x_{\text{full}} \leftarrow x + h\,f_{\text{drift}}(x)$
\State $x_{\text{half}} \leftarrow x + (h/2)\,f_{\text{drift}}(x)$
\State compute drift $f_{\text{drift}}(x_{\text{half}})$
\State $\tilde{x}_{\text{full}} \leftarrow x_{\text{half}} + (h/2)\,f_{\text{drift}}(x_{\text{half}})$
\State estimate relative error $\rho = \|\tilde{x}_{\text{full}} - x_{\text{full}}\| / (\|x\| + \epsilon)$
\State $m \leftarrow \mathrm{clip}(\lceil \sqrt{\rho/\text{Tol}} \rceil, m_{\min}, m_{\max})$
\If{$m \le 2$}
  \State $x \leftarrow \tilde{x}_{\text{full}}$ \Comment{reuse current predictor if accurate}
\Else
  \For{$s=1$ to $m$}
    \State compute drift $f_{\text{drift}}(x)$
    \State $x \leftarrow x + (h/m)\,f_{\text{drift}}(x)$
  \EndFor
\EndIf
\State \Return $x$
\end{algorithmic}
\end{algorithm}
